\chapter{Incidence dodécaphasique et corridor Witt--Leech--Monstre}
\label{chap:bandera-excepcional-acumulativa}
\index{double projection!drapeau commun}
\index{Witt!plan}
\index{Monstre}

Le \cref{chap:paley-codigo-witt} a construit explicitement la matrice
\(A_W\), le code \(C_W\) et le système \(W_{12}\). Sur ces coordonnées,
la double projection se formalise au moyen de deux applications d'évaluation qui,
dans le calibre déclaré, produisent par des opérations distinctes le même drapeau
polarisé de Witt. Ce drapeau calibré détermine en outre une origine marquée. La
construction atteint le plan \(W_{12}\), son groupe d'automorphismes
\(M_{12}\) et, par une seconde branche qui ne passe pas par \(W_{24}\), un
3-voisin pair unimodulaire sans racines. La classification classique identifie ce
dernier au réseau de Leech ; l'extension à l'algèbre d'opérateurs de
vertex lunaire et au groupe Monstre s'obtient au moyen des théorèmes classiques
correspondants.

\section{État dodécaphasique et supports d'incidence}

Le vecteur dodécaphasique est
\[
K=(234,543,140,729,659,824,621,058,914,794,146,601).
\]
Les quatre mots du cadre sont
\[
w_\pi=010211,
\quad
w_e=201101,
\quad
w_\varphi=121200,
\quad
w_{\rm APP}=022020.
\]
Dans le calibre symétrique de Paley--Witt, on utilise l'encodeur
\[
c(w)=(w,wA_W)\in\mathbb F_3^{12}.
\]
Leurs supports sont
\[
\begin{aligned}
P_\pi&=\{2,4,5,6,7,8,9,10,11\},\\
H_e&=\{1,3,4,6,9,10\},\\
H_\varphi&=\{1,2,3,4,7,9\},\\
H_{\rm APP}&=\{2,3,5,10,11,12\}.
\end{aligned}
\]
\(P_\pi\) a neuf points. Les trois autres supports sont des hexades. Cette
distinction doit être conservée : les douze hexades contenues dans \(P_\pi\) ont
pour triplets complémentaires les droites de \(AG(2,3)\).

\section{Coïncidence des supports de seuil}

L'évaluation archimédienne extrait du vecteur \(K\) le support de seuil
\[
\beta_{\rm arq}(K)=\{i:K_i\ge729\}
=\{4,6,9,10\}.
\]
L'évaluation exceptionnelle calcule l'intersection
\[
\beta_{\rm exc}(w_e,w_\pi)
=\operatorname{supp}c(w_e)\cap\operatorname{supp}c(w_\pi).
\]
L'évaluation produit la première égalité exacte conditionnée par ces données :
\begin{equation}
\boxed{
\beta_{\rm arq}(K)
=\beta_{\rm exc}(w_e,w_\pi)
=B_K=\{4,6,9,10\}.}
\label{eq:bandera-cara-alta}
\end{equation}
Les deux applications réalisent des opérations distinctes sur des entrées
déclarées : la première utilise des seuils du vecteur \(K\), tandis que la seconde
utilise des intersections de supports du code. Le seuil \(729\), les mots et
le calibre de Paley font partie de la calibration. En particulier, cette
égalité n'est pas une validation indépendante de \(K\) : tous les seuils
entiers entre \(660\) et \(729\), inclus, sélectionnent le même support
\(B_K\).

\section{Coïncidence des supports négatifs}

Avant la retenue, la cochaîne archimédienne est
\[
s^{(0)}=P+E-\Phi-K,
\]
et ses douze composantes sont
\[
(7,297,353,-430,-715,-1199,-3,286,-894,-528,380,663).
\]
Son support négatif vaut
\[
\eta_{\rm arq}(s^{(0)})
=\{4,5,6,7,9,10\}.
\]
Par la voie exceptionnelle, nous fixons comme sélecteur de calibration l'hexade \(H\)
satisfaisant
\[
B_K\subset H\subset P_\pi
\]
et
\[
(|H\cap H_e|,|H\cap H_\varphi|,|H\cap H_{\rm APP}|)=(4,3,2).
\]
L'énumération du plan produit une solution unique pour ce sélecteur :
\[
\eta_{\rm exc}
=\{4,5,6,7,9,10\}
=H^-_\alpha.
\]
Par conséquent,
\begin{equation}
\boxed{
\eta_{\rm arq}(s^{(0)})
=\eta_{\rm exc}(B_K,P_\pi,H_e,H_\varphi,H_{\rm APP})
=H^-_\alpha.}
\label{eq:bandera-hexada-negativa}
\end{equation}

Les équations~\eqref{eq:bandera-cara-alta} et
\eqref{eq:bandera-hexada-negativa} se condensent dans le drapeau
\[
\mathcal I_*^{\rm core}=(B_K\subset H^-_\alpha).
\]
C'est le drapeau de compatibilité calibrée de la double projection. Le cadre marqué
\[
\widetilde{\mathcal I}_*
=(\mathcal I_*^{\rm core};P_\pi,H_e,H_\varphi,H_{\rm APP})
\]
conserve les données nécessaires pour construire la branche exceptionnelle. Sous une
permutation simultanée \(g\in M_{12}\), toute la construction est transportée
de manière équivariante ; l'invariant est sa classe d'isomorphisme, non les noms
numériques isolés.

\section{Reconstruction du point de base à partir du cadre d'incidence}

L'origine n'a pas à rester un sélecteur extérieur si le cadre
ordonné complet est conservé. Nous définissons
\[
o(\mathcal F)
=(H_e\cap H_\varphi)\setminus(H^-\cup H_{\rm APP}),
\qquad
\mathcal F=(H^-,H_e,H_\varphi,H_{\rm APP}).
\]

\begin{theorem}[Origine déterminée par les incidences]
Pour le cadre canonique,
\[
o(\mathcal F_*)=\{1\}.
\]
De plus, la règle est équivariante : \(o(g\mathcal F)=g\,o(\mathcal F)\).
\end{theorem}

\begin{proof}
On a
\[
H_e\cap H_\varphi=\{1,3,4,9\},
\]
et
\[
(H^-_\alpha\cup H_{\rm APP})\cap\{1,3,4,9\}
=\{3,4,9\}.
\]
La différence est \(\{1\}\). L'équivariance découle du fait que l'intersection,
l'union et la différence commutent avec l'action de toute permutation.
\end{proof}

Le passage de cette origine d'incidence à un vecteur d'un réseau
\(A_2^{12}\) exige de fixer une réalisation métrique des douze coordonnées.
Le \cref{sec:vecino-leech-ejecutable} fixe cette réalisation, montre
que l'origine sélectionne l'un des douze minimiseurs radiaux et certifie
les conditions du 3-voisin. L'incidence détermine l'origine ; la matrice de
Gram et la chambre positive déterminent ensuite le vecteur.

\section{Incidence du plan ternaire}

La chaîne de Paley
\[
C_P\longmapsto A_W\longmapsto C_W\longmapsto W_{12}
\]
est fixée par les matrices et démonstrations du
\cref{chap:paley-codigo-witt}. En particulier,
\(A_W^2=-I_6\), \(C_W=[12,6,6]_3\), et les 132 supports projectifs de
poids six forment \(W_{12}=S(5,6,12)\). Le passage d'un mot à son support
identifie exactement la paire \(\{c,-c\}\) ; les incidences utilisées dans ce
qui suit appartiennent donc au plan projectif et non à des mots
signés individuels.

\section{Hexades, octades et identité de cardinalité 28}

Le plan binaire \(W_{24}=S(5,8,24)\) possède des octades, chacune comportant
\(\binom82=28\) paires. Le
\cref{p12-83-c17-s1:thm:dodecada-w12} construit
\(W_{12}\) comme plan résiduel d'une dodécade \(D\) du code binaire
étendu de Golay ; le \cref{thm:elevacion-unica-hexada} démontre qu'une
fois fixé l'alignement \(\ell\), chaque hexade possède un relèvement unique
\(O_H\) à une octade. Le passage \(6\to8\) est donc un morphisme
d'incidence relatif à la donnée binaire \((D,\ell)\), et non une inférence fondée
sur les cardinalités.

D'autre part, la microfibre de \(\pi\) contient 1008 voies et une nonade a
36 paires, de sorte que
\[
1008=36\cdot28.
\]
L'identité \(1008=36\cdot28\) est un contrôle numérique du catalogue. Le
morphisme d'incidence est défini indépendamment au moyen de
\(H\mapsto\iota_D(\ell(H))\) ; l'identité des cardinalités n'intervient pas dans
sa définition.

\section{Involutions d'étoile et génération du groupe de Mathieu}

Pour chaque face \(B\in\binom{[12]}4\), il existe quatre hexades de \(W_{12}\)
contenant \(B\). Les quatre résidus \(H\setminus B\) sont des paires disjointes
qui partitionnent \([12]\setminus B\). L'involution \(\sigma_B\) fixe \(B\)
et échange les extrémités de ces quatre paires. Il y a
\[
\binom{12}{4}=495
\]
involutions d'étoile.

L'énumération exhaustive, pour la matrice et les mots déclarés, vérifie
\[
\sigma_B(W_{12})=W_{12}
\]
pour toutes et
\[
\left|\langle\sigma_B:B\in\tbinom{[12]}4\rangle\right|=95040.
\]
Comme \(\Aut(W_{12})\cong M_{12}\), les 495 étoiles engendrent \(M_{12}\).
La face HMT est
\[
B_K=\{4,6,9,10\},
\qquad
\sigma_{B_K}=(1\ 3)(2\ 11)(5\ 7)(8\ 12).
\]
Une seule étoile n'engendre que \(C_2\). L'action combinatoire globale
s'exprime au moyen de la représentation du groupe libre des générateurs d'étoile
\[
\rho_\star:F_{495}\longrightarrow\Aut(W_{12}),
\qquad
\rho_\star(\gamma_B)=\sigma_B,
\qquad
\operatorname{im}\rho_\star=M_{12},
\]
où \(F_{495}\) est le groupe libre sur les 495 étoiles formelles. Le groupe
global est donc l'image de cette famille d'involutions et non celle d'une
seule étoile. Identifier ces générateurs à des lacets de l'espace des états
exigerait de construire en outre le transport correspondant ; cette
identification n'intervient pas dans le calcul du groupe.

\section{Deux prolongements exceptionnels du code ternaire}

À partir du code ternaire, deux branches de types distincts peuvent être prolongées :
\[
\begin{aligned}
C_W&\longrightarrow W_{12}
\xrightarrow[\ (D,\ell)\ ]{\ \iota_D\ }W_{24},\\
C_W&\xrightarrow[\phi]{\text{recollement explicite}}N(A_2^{12})
\xrightarrow[v_\alpha]{\text{3-voisin exact}}\Lambda_{24}
\longrightarrow V^\natural .
\end{aligned}
\]
La première branche se poursuit au moyen du plan résiduel de la dodécade et de son
relèvement unique ; la seconde se poursuit au moyen d'un recollement ternaire et d'un voisin
d'indice trois. Les deux constructions ont des codomaines distincts. Pour la
seconde, soit \(A_2=\mathbb Z\alpha_1\oplus\mathbb Z\alpha_2\) de matrice
de Gram
\[
G_{A_2}=\begin{pmatrix}2&-1\\-1&2\end{pmatrix},
\qquad \rho=\alpha_1+\alpha_2,
\]
et prenons comme représentants de \(A_2^*/A_2\cong\mathbb F_3\)
\[
\varpi_0=0,\qquad
\varpi_1=\frac{2\alpha_1+\alpha_2}{3},\qquad
\varpi_2=\frac{\alpha_1+2\alpha_2}{3}.
\]
La réalisation du recollement est l'application coordonnée par coordonnée
\[
\phi(c_1,\ldots,c_{12})=(\varpi_{c_1},\ldots,\varpi_{c_{12}}),
\]
qui forme
\[
N(C_W)=\bigcup_{c\in C_W}(A_2^{12}+\phi(c)).
\]
La linéarité de \(C_W\) fait de l'union des classes un réseau ; son
autodualité donne l'intégralité du recollement, et la divisibilité par trois de
tous les poids donne sa parité. De plus, \(C_W\) est de dimension six et
\[
\det N(C_W)=\frac{3^{12}}{|C_W|^2}
=\frac{3^{12}}{(3^6)^2}=1.
\]
La distance minimale six implique qu'une classe de recollement non nulle a une norme
au moins égale à \(6(2/3)=4\). Les racines de \(N(C_W)\) sont donc exactement celles
de \(A_2^{12}\). C'est la construction explicite du Niemeier de type
\(A_2^{12}\) utilisée dans la suite.

\section{Minimisation radiale et construction explicite du 3-voisin}
\label{sec:vecino-leech-ejecutable}

Dans la chambre positive, nous considérons des vecteurs radiaux
\[
v(a)=\sum_{i=1}^{12}a_i\rho_i,
\qquad a_i=1+3b_i,\quad b_i\ge0.
\]
La condition \(v(a)^2\equiv0\pmod{18}\) équivaut à
\(\sum_i b_i\equiv1\pmod3\). Son minimum est
\[
v(a)^2=2\sum_i(1+3b_i)^2=54,
\]
et est atteint exactement douze fois : un coefficient vaut quatre et les onze
restants valent un. L'origine retrouvée \(o_*=1\) sélectionne
\[
v_\alpha=4\rho_1+\rho_2+\cdots+\rho_{12}.
\]
La minimalité appartient au cône radial positif déclaré. Dans tout le
quotient résiduel, il existe, par exemple, le représentant
\((\alpha_1-2\alpha_2,\rho,\ldots,\rho)\), congru à
\((\rho,\ldots,\rho)\) modulo trois et de norme \(14+11\cdot2=36\). Le
calcul direct le conserve comme témoin qui empêche d'élargir le domaine.

On définit
\[
N_{\alpha,0}=\{x\in N(C_W):
\langle x,v_\alpha\rangle\equiv0\pmod3\},
\qquad
N_\alpha=N_{\alpha,0}+\mathbb Z\frac{v_\alpha}{3}.
\]

\begin{theorem}[3-voisin de Leech]
La construction précédente satisfait
\begin{align*}
\mathrm{(I1)}\quad&v_\alpha\in A_2^{12}\subset N(C_W),
&\langle N(C_W),v_\alpha\rangle&\not\subset3\mathbb Z,\\
\mathrm{(I2)}\quad&v_\alpha^2=54\equiv0\pmod{18},\\
\mathrm{(I3)}\quad&v_\alpha/3\in
N_{\alpha,0}^*\setminus N(C_W),
&3(v_\alpha/3)&\in N_{\alpha,0}.
\end{align*}
Le réseau \(N_\alpha\) est pair, unimodulaire, de rang vingt-quatre et sans
racines. Par conséquent,
\[
N_\alpha\cong\Lambda_{24}.
\]
\end{theorem}

\begin{proof}
Pour toute racine d'un facteur \(A_2\), l'accouplement avec
\(a_i\rho_i\) est congru à \(1\) ou \(2\) modulo trois. Le caractère de
(I1) est donc surjectif, et aucune des anciennes racines n'appartient à
\(N_{\alpha,0}\). L'égalité de (I2) est le calcul radial précédent. Si
\(x\in N_{\alpha,0}\), alors
\(\langle x,v_\alpha/3\rangle\in\mathbb Z\), ce qui démontre la première
inclusion de (I3). Si \(v_\alpha/3\) appartenait à \(N(C_W)\),
l'intégralité du Niemeier imposerait la nullité du caractère de (I1). La
classe est d'ordre trois et \(v_\alpha\in N_{\alpha,0}\) car
\(v_\alpha^2\equiv0\pmod3\).

Le noyau est d'indice trois, donc
\[
\det N_{\alpha,0}=3^2\det N(C_W)=9,
\qquad
\det N_\alpha=9/3^2=1.
\]
Pour \(x\in N_{\alpha,0}\) et \(m\in\mathbb Z\),
\[
\left(x+m\frac{v_\alpha}{3}\right)^2
=x^2+2m\frac{\langle x,v_\alpha\rangle}{3}+6m^2
\in2\mathbb Z,
\]
de sorte que le voisin est pair. Il reste à exclure de nouvelles racines. Chaque coordonnée des
deux classes non triviales appartient à l'une des six classes
\[
A_2+\varpi_j\pm\rho/3,\qquad j\in\{0,1,2\},
\]
et le minimum de chaque classe est \(2/9\). En effet, en écrivant une coordonnée
sous la forme \((u/3,v/3)\), la norme est
\(2(u^2-uv+v^2)/9\) ; les six résidus sont non nuls et chacun admet un
représentant avec \(u^2-uv+v^2=1\). Toute
nouvelle norme est au moins \(12(2/9)=8/3>2\). Le voisin n'a pas de racines. La
classification des réseaux pairs unimodulaires de rang vingt-quatre
l'identifie à Leech.
\end{proof}

L'énumération des six minima locaux et des douze minimiseurs radiaux
reconstruit les \(729\) mots de \(C_W\), la matrice de Gram du recollement et
sa parité ; l'arithmétique exacte vérifie I1--I3, les déterminants et la borne
\(8/3\).

\section{Corridor de Leech à l'algèbre Moonshine}
\label{sec:corredor-hmt-moonshine}

Définissons
\[
 \Lambda_{\rm HMT}:=N_\alpha\cong\Lambda_{24},
 \qquad
 V^\natural_{\rm HMT}
 :=V_{\Lambda_{\rm HMT}}^+
 \oplus V_{\Lambda_{\rm HMT}}^{T,+}.
\]
Aucun second réseau n'est introduit ici : la construction
Frenkel--Lepowsky--Meurman s'applique au voisin produit dans la sous-section
précédente.

\begin{theorem}[Corridor HMT--Moonshine]
\label{thm:corredor-hmt-moonshine}
Une fois fixés la dodécade, l'alignement, l'origine et la chambre exceptionnelle de la
construction précédente, on a
\[
 V^\natural_{\rm HMT}\cong V^\natural,
 \qquad
 \Aut(V^\natural_{\rm HMT})\cong\mathbb M,
\]
et
\[
 \operatorname{Tr}_{V^\natural_{\rm HMT}}q^{L_0-1}
 =J_{\rm Mon}(\tau)
 =j(\tau)-744
 =q^{-1}+196884q+21493760q^2+\cdots .
\]
Par conséquent, la chaîne
\[
 \mathrm{HMT}\longrightarrow N_\alpha\cong\Lambda_{24}
 \longrightarrow V^\natural_{\rm HMT}
 \longrightarrow\mathbb M
 \longrightarrow\text{\rm Moonshine}
\]
est un résultat démontré par composition.
\end{theorem}

\begin{proof}
Le théorème du 3-voisin démontre, à partir des données déclarées, que
\(N_\alpha\) est pair, unimodulaire, de rang vingt-quatre et sans racines. La
classification des réseaux l'identifie à \(\Lambda_{24}\). Appliquée à
cette sortie, la construction de Frenkel--Lepowsky--Meurman produit
l'orbifold
\(V_{\Lambda_{\rm HMT}}^+\oplus V_{\Lambda_{\rm HMT}}^{T,+}\) ;
les théorèmes classiques identifient son groupe d'automorphismes à
\(\mathbb M\) et son caractère gradué au Hauptmodul normalisé
\(J_{\rm Mon}=j-744\)
\cite{FLM1988,ConwayNorton1979,Borcherds1992}. Chaque flèche fournit donc
un objet du type requis par la suivante.\qedhere
\end{proof}

\begin{remark}[La composante triviale n'est pas un ajustement]
La graduation satisfait
\[
 V^\natural_{\rm HMT}
 =\bigoplus_{n\geq0}(V^\natural_{\rm HMT})_n,
 \qquad
 (V^\natural_{\rm HMT})_0=\CC\mathbf1,
 \qquad
 (V^\natural_{\rm HMT})_1=0.
\]
Le terme \(q^{-1}\) enregistre ainsi la droite du vide. En degré deux,
\[
 \dim(V^\natural_{\rm HMT})_2
 =196884
 =1+196883
 =196560+324
 =54(5\cdot729+1).
\]
L'égalité \(196884=1+196883\) sépare la représentation triviale et la plus petite
représentation irréductible non triviale du Monstre. Les deux autres égalités
sont des décompositions entières exactes ; elles ne remplacent pas cette structure graduée
et ne sont pas utilisées pour sélectionner rétrospectivement le corridor.
\end{remark}

L'étape finie jusqu'à \(N_\alpha\) part d'une dodécade, d'un alignement, d'une
origine et d'une chambre déclarés. Le calcul prouve que \(N_\alpha\) est pair,
unimodulaire et sans racines. Les théorèmes classiques de
classification et d'orbifold s'appliquent à cette sortie : l'orbifold est \(V^\natural\), son groupe
d'automorphismes est \(\mathbb M\) et son caractère gradué est \(J_{\rm Mon}\).
La composition conserve ainsi la provenance de chaque flèche sans réduire le
corridor à une coïncidence de coefficients.

\section{Espace quotient de 243 syndromes}

Poncturer une coordonnée de \(C_W\) produit
\[
G_{11}=[11,6,5]_3.
\]
La boule ternaire de rayon deux a
\[
1+22+220=243
\]
éléments. Comme
\[
|G_{11}|\cdot243=3^6\cdot3^5=3^{11},
\]
et la distance minimale est cinq, ces boules sont disjointes et couvrent tout
\(\mathbb F_3^{11}\). Chaque vecteur admet une décomposition unique
\[
x=c+e,
\qquad
c\in G_{11},
\quad
\operatorname{wt}(e)\le2.
\]
Le quotient linéaire est
\[
\mathbb F_3^{11}/G_{11}\cong\mathbb F_3^5,
\qquad
|\mathbb F_3^{11}/G_{11}|=243.
\]
On obtient donc un espace quotient de 243 syndromes dans un espace ambiant de onze
coordonnées et de dimension quotient cinq. Le catalogue de 243 mots \(w_6\)
et ce quotient ont le même cardinal, mais sont des objets de types distincts ;
l'encodeur direct envoie les mots du code sur le syndrome nul. La double projection
démontrée utilise le drapeau commun, non une bijection de cardinalités sans
structure.

\section{Compatibilité finie des évaluations arithmétique et exceptionnelle}

\begin{theorem}[Compatibilité finie des deux projections]
Considérons le vecteur \(K\) et l'état dodécaphasique signé construits dans
le \cref{chap:alpha} ; les quatre mots du catalogue définis dans le
\cref{chap:regiones} ; et la matrice \(A_W\) construite dans le
\cref{chap:paley-codigo-witt}. La construction fixe une seule calibration
\(c_{\rm exc}\in\mathscr C_{\rm exc}\) et la transporte avec l'état ; cette
calibration type les coordonnées du codomaine exceptionnel et n'agit pas comme
entrée génératrice. Les énoncés suivants sont satisfaits :
\begin{enumerate}
  \item l'évaluation archimédienne obtient \(\alpha_{12}\) au moyen de la retenue
  dodécaphasique ;
  \item les applications de support archimédien et exceptionnel produisent le même drapeau
  \(B_K\subset H^-_\alpha\) ;
  \item le cadre marqué retrouve l'origine \(o_*=1\) ;
  \item la matrice \(A_W\) construit \(W_{12}\), et les 495 involutions
  d'étoile engendrent un groupe d'ordre \(95040\), identifié à \(M_{12}\) ;
  \item la réalisation discriminante explicite construit \(N(A_2^{12})\),
  et l'origine, avec la chambre positive, sélectionne \(v_\alpha\) ; son
  3-voisin est pair, unimodulaire et sans racines.
\end{enumerate}
Les identités des supports et la règle de l'origine sont équivariantes sous une
permutation simultanée des douze coordonnées. La calibration transforme de
manière simultanée les données d'incidence et ne modifie pas leur classe d'isomorphisme.
\end{theorem}

\begin{proof}
Le premier énoncé est la reconstruction dodécaphasique du
\cref{chap:alpha}. Les équations
\eqref{eq:bandera-cara-alta} et \eqref{eq:bandera-hexada-negativa} démontrent
le deuxième par évaluation directe des deux applications. Le calcul ensembliste
de la section précédente donne \(o_*=1\). Enfin, l'énumérateur, les 132
supports de poids six, les 495 involutions et l'ordre du groupe, avec le
théorème du 3-voisin, prouvent le dernier énoncé.
\end{proof}

\begin{corollary}[Compatibilité avec le corridor HMT--Moonshine]
Les données finies de la projection exceptionnelle construisent le même
\(N_\alpha\) qui intervient dans le
\cref{thm:corredor-hmt-moonshine} ; par conséquent, l'évaluation finie est
compatible avec \(V^\natural_{\rm HMT}\), son action de \(\mathbb M\) et son
caractère \(J_{\rm Mon}\).
\end{corollary}


Le drapeau fini s'obtient au moyen de deux opérations distinctes appliquées aux
coordonnées du même état enrichi, et l'origine est dérivée après
fixation de ce drapeau. Cela prouve une identité interne typée et conserve la
provenance commune sans élever la calibration au rang d'entrée causale. La
construction réticulaire conserve la différence de type. La composition depuis
les données de \(\alpha\) jusqu'à \(\Lambda_{24}\) se factorise par le code, le
recollement \(A_2^{12}\), l'origine marquée, la chambre radiale et le 3-voisin. La
branche \(W_{12}\to W_{24}\) se poursuit séparément et est déterminée pour tout
choix d'une dodécade \(D\) et d'un isomorphisme d'incidence
\(\ell:W_{12}^{\rm HMT}\to\mathcal H(D)\). L'équivariance démontrée fait
que ces choix coordonnent la représentation sans modifier la classe
d'isomorphisme engendrée ni la fermeture de la double projection.
